% from ega4-2-fr.tex
\oldpage[IV-2]{140}

Si de plus $M\neq0$ et $N\neq0$, on a
\begin{equation}
\operatorname{coprof}_B(M\otimes_AN)=\operatorname{coprof}_A(M).
\tag{6.3.4.2}\label{IV.2.s6.3.4.2-fr}
\end{equation}
\end{corollary}

Il revient en effet au même de dire que $N\otimes_Ak$ est un $B$-module de
longueur finie ou un $(B\otimes_Ak)$-module de longueur finie, et l'on sait
(0, 16.2.3) que les modules de longueur finie $\neq0$ sont de dimension $0$ et
de profondeur $0$.

Le corollaire (6.3.4) sera applicable en particulier lorsque $B\otimes_Ak$ est
un anneau de longueur finie, c'est-à-dire lorsque $\mathfrak mB$ est un idéal
de définition de l'anneau $B$.

\begin{corollary}[6.3.5]
\label{IV.2.s6.3.5-fr}
Soient $X$, $Y$ deux préschémas localement noethériens, $f:X\to Y$ un
morphisme plat.
\begin{enumerate}
\item[(i)] Si, en un point $x\in X$, on a
$\operatorname{coprof}(\mathcal O_x)\leq n$, alors, en posant $y=f(x)$, on a
$\operatorname{coprof}(\mathcal O_y)\leq n$ et
$\operatorname{coprof}(\mathcal O_x\otimes_{\mathcal O_y}k(y))\leq n$; en
particulier, si $\mathcal O_x$ est un anneau de Cohen-Macaulay, il en est de
même de $\mathcal O_y$ et de
$\mathcal O_x\otimes_{\mathcal O_y}k(y)$.
\item[(ii)] Supposons inversement que
$\mathcal O_x\otimes_{\mathcal O_y}k(y)$ soit un anneau de Cohen-Macaulay.
Alors, si $\operatorname{coprof}(\mathcal O_y)\leq n$ (resp. si
$\mathcal O_y$ est un anneau de Cohen-Macaulay), on a
$\operatorname{coprof}(\mathcal O_x)\leq n$ (resp. $\mathcal O_x$ est un
anneau de Cohen-Macaulay).
\end{enumerate}
\end{corollary}

Il suffit d'appliquer (6.3.2) pour $M=A=\mathcal O_y$ et
$N=B=\mathcal O_x$.

\begin{corollary}[6.3.6]
\label{IV.2.s6.3.6-fr}
Soit $A$ un anneau de Cohen-Macaulay. Alors
$A'=A[T_1,\ldots,T_n]$ ($T_i$ indéterminées) est un anneau de
Cohen-Macaulay.
\end{corollary}

En effet, si l'on pose $Y=\operatorname{Spec}(A)$,
$X=\operatorname{Spec}(A')$ et si $f:X\to Y$ est le morphisme canonique, $f$
est plat (puisque $A'$ est un $A$-module libre (2.1.2)) et pour tout $y\in Y$,
$k(y)[T_1,\ldots,T_n]$ est un anneau régulier (0, 17.3.7), donc \emph{a
fortiori} de Cohen-Macaulay; il suffit donc d'appliquer (6.3.5).

\begin{corollary}[6.3.7]
\label{IV.2.s6.3.7-fr}
Tout quotient d'un anneau de Cohen-Macaulay est universellement caténaire.
\end{corollary}

Cela résulte de (6.3.6), de (5.6.1) et de (0, 16.5.12).

\begin{proposition}[6.3.8]
\label{IV.2.s6.3.8-fr}
Soit $A$ un anneau local noethérien; supposons qu'il existe un $A$-module $M$
de type fini qui soit un $A$-module de Cohen-Macaulay et dont le support soit
égal à $\operatorname{Spec}(A)$. Soit $B$ un anneau quotient de $A$, et soit
$f:\operatorname{Spec}(\widehat B)\to\operatorname{Spec}(B)$ le morphisme
canonique; alors, pour tout $x\in\operatorname{Spec}(B)$, $f^{-1}(x)$ est un
schéma de Cohen-Macaulay.
\end{proposition}

Comme $B$ est un $A$-module de type fini, on a
$\widehat B=B\otimes_A\widehat A$ $(0_{\mathrm I}, 7.3.3)$, donc $f$ est la
restriction à $\operatorname{Spec}(\widehat B)$ du morphisme canonique
$\operatorname{Spec}(\widehat A)\to\operatorname{Spec}(A)$, et les fibres de
ces deux morphismes en un point de $\operatorname{Spec}(B)$ sont donc les
mêmes. On est par suite ramené à démontrer la proposition lorsque $B=A$.
Soit donc $\mathfrak p=\mathfrak j_x$ un idéal premier de $A$; comme par
hypothèse $\mathfrak p\in\operatorname{Supp}(M)$, on sait (0, 16.5.6 et
16.5.9) qu'il existe une suite $M$-régulière $(t_i)_{1\leq i\leq r}$
d'éléments de $\mathfrak p$ telle que
\[
N=M/\Bigl(\sum_{i=1}^{r}t_iM\Bigr)
\]
soit un $A$-module de Cohen-Macaulay, $\mathfrak p$ un idéal premier associé
minimal de $N$ et
$\dim(N)=\dim(M/\mathfrak pM)=\dim(A/\mathfrak p)$. Le même raisonnement
qu'au début montre que l'on peut remplacer $M$ par $N$ et $A$ par
$A/(\sum_{i=1}^{r}t_iA)$, et par suite supposer que $\mathfrak p$ est un

